The last piece of the puzzle to be added to complete the approach specification is the implementation of the 
\textit{dummy shared memory}. This is a temporary memory region used mainly to prevent the program from crashing 
when certain circumstances are met.\\
In particular, especially in stripped binaries, it is difficult to correctly identify which functions are the actual 
program code and which functions are used by default to make the program run (e.g. \textit{\_start, frame\_dummy \dots}).
For this reason, it sometimes happens that some functions that are executed prior to the main are actually instrumented 
by the tool.
Since each \gls{IP} holds the instructions that perform the \gls{SHM} update according to \cref{alg:shm_update}, 
at each \gls{IP} there is an access to the \textit{shared\_map\_address} saved into the created section.
However, as seen in \cref{sec:actual_instrumentation}, that address is actually written only when the Initial Instrumentation 
is reached at the beginning of the main.\\
The idea is then to create a new empty section and write the address of this memory region in the location where 
the shared memory address should be written.
In this way, if at any point prior to the main some instruction tries to access the shared memory address, it will 
actually access the address of the dummy shared memory and update a memory region that is not used, so the binary 
does not crash.
At the beginning of the main, the actual address of the \gls{SHM} is dereferenced, copied from the \gls{GOT}, and 
overwrites the address of the dummy shared memory. From this moment on, each \gls{BB} updates the real \gls{SHM}.\\
Depending on whether the binary is \gls{PIE} or not, the address of the dummy section must be written differently:
\begin{itemize}
\item \textbf{\gls{PIE}}: in the presence of \gls{ASLR}, the address at which the binary will be loaded is not known a priori, 
            and so it is impossible to know in advance the address at which the dummy section will be loaded.
            For this reason, a relocation must be added in order to write the correct address at load time.
\item \textbf{non-\gls{PIE}}: in this case the correct address is known at link time, so it is possible to write 
        the address of the dummy section directly in the location where the \gls{SHM} address will be written.
\end{itemize}